\PassOptionsToPackage{dvipsnames}{xcolor}

\documentclass[12pt]{beamer}
\usetheme{default}
\usecolortheme{crane}

\usepackage[T1]{fontenc}
\usepackage[slovak]{babel}

\usepackage{amsmath}
\usepackage{amsmath, amssymb}
\usepackage{hyperref, url}
\usepackage{graphicx}
\usepackage{array}
\usepackage{alltt}

%\setbeamersize{text margin left=1pt,text margin right=1pt}
\setbeamertemplate{footline}[frame number]
\beamertemplatenavigationsymbolsempty

% https://www.overleaf.com/learn/latex/Using_colours_in_LaTeX
\def\blue#1{\textcolor{Cerulean}{#1}}
\def\green#1{\textcolor{LimeGreen}{#1}}

% database-related stuff
\DeclareMathOperator{\join}{\bowtie}
\DeclareMathOperator{\antijoin}{\rhd}

\DeclareMathOperator{\lubi}{lubi}
\DeclareMathOperator{\capuje}{capuje}
\DeclareMathOperator{\navstivil}{navstivil}
\DeclareMathOperator{\vypil}{vypil}
\DeclareMathOperator{\answer}{answer}

\title{Window functions, polia, JSON}
\author{Ján Mazák}
\institute{FMFI UK Bratislava}
\date{}


\begin{document}

\frame{\titlepage}

\begin{frame}[fragile]{Window functions}
\begin{itemize}
\item Použitie GROUP BY vedie k nahradeniu celej skupiny jediným riadkom.
\item Niekedy však chceme zachovať pôvodné riadky a len k nim doplniť hodnotu súvisiacu s inými riadkami v skupine.
\item Riešenie: \alert{window functions}. Tieto funkcie vidia nielen jeden riadok, ale aj \uv{okolité} riadky
    (presná definícia je súčasťou dotazu v časti \alert{OVER}).
\item Window functions nevidia riadky odfiltrované vo WHERE a HAVING.
\item Ak chceme filtrovať výsledok podľa vypočítanej hodnoty window funkcie, treba použiť vnorený dotaz.
\end{itemize}
\end{frame}

\begin{frame}[fragile]{Window functions}
\small
\begin{alltt}
SELECT emp_id, name, dept_id, salary,
       \green{avg}(salary) \alert{OVER (PARTITION BY dept_id)} avg_s
FROM employee;
\end{alltt}
\begin{alltt}
\blue{
 emp_id | name   | dept_id | salary  | avg_s
-{}-{}-{}-{}-{}-{}-{}-{}+-{}-{}-{}-{}-{}-{}-{}-{}+-{}-{}-{}-{}-{}-{}-{}-{}-{}+-{}-{}-{}-{}-{}-{}-{}-{}-{}+-{}-{}-{}-{}-{}-{}
     11 | King   |      10 | 5000.00 | 2600.00
     13 | Clark  |      10 | 1500.00 | 2600.00
     24 | Miller |      10 | 1300.00 | 2600.00
     14 | Jones  |      20 | 2975.00 | 2175.00
     ...
}
\end{alltt}
\end{frame}

\begin{frame}[fragile]{Window functions}
\footnotesize
\begin{alltt}
SELECT dept_id, name, salary,
       \green{rank()} \alert{OVER (PARTITION BY dept_id
                    ORDER BY salary DESC)}
FROM employee;
\end{alltt}
\begin{alltt}
\blue{
 dept_id |  name  | salary  | rank
-{}-{}-{}-{}-{}-{}-{}-{}-{}+-{}-{}-{}-{}-{}-{}-{}-{}+-{}-{}-{}-{}-{}-{}-{}-{}-{}+-{}-{}-{}-{}-{}
      10 | King   | 5000.00 |    1
      10 | Clark  | 1500.00 |    2
      10 | Miller | 1300.00 |    3
      20 | Ford   | 3000.00 |    1
      20 | Scott  | 3000.00 |    1
      20 | Jones  | 2975.00 |    3
      20 | Adams  | 1100.00 |    4
    ...
}
\end{alltt}
\end{frame}

\begin{frame}[fragile]{Window functions: dense\_rank(), row\_number()}
\footnotesize
\begin{alltt}
SELECT name, salary, rank() OVER \blue{w},
       \green{dense_rank()} OVER \blue{w}, \green{row_number()} OVER \blue{w}
FROM employee
WINDOW \blue{w} AS (PARTITION BY dept_id ORDER BY salary DESC);
\end{alltt}
\begin{alltt}
\blue{
  name  | salary  | rank | dense_rank | row_number
-{}-{}-{}-{}-{}-{}-{}-{}+-{}-{}-{}-{}-{}-{}-{}-{}-{}+-{}-{}-{}-{}-{}-{}+-{}-{}-{}-{}-{}-{}-{}-{}-{}-{}-{}-{}+-{}-{}-{}-{}-{}-{}-{}-{}-{}-{}-{}
 Ford   | 3000.00 |    1 |          1 |          1
 Scott  | 3000.00 |    1 |          1 |          \alert{2}
 Jones  | 2975.00 |    3 |          \alert{2} |          3
 Adams  | 1100.00 |    4 |          3 |          4
 ...
}
\end{alltt}
\scriptsize
\begin{itemize}
\item dense\_rank(): po zhode nevynecháva čísla.
\item row\_number(): čísluje 1, 2, 3, \dots{} aj pri zhode; poradie zhodných riadkov
  je ľubovoľné, deterministicky: ORDER BY salary DESC, emp\_id.
\end{itemize}
\end{frame}

\begin{frame}[fragile]{Window functions}
Funkcie, ktoré možno použiť ako window functions:
\begin{itemize}
\item bežné agregačné funkcie\\
    (môžu sa však správať trochu odlišne!)
\item rank(), dense\_rank(), row\_number(), percent\_rank()
\item nth\_value() --- element in the n-th row of the window
\item \dots
\end{itemize}
\end{frame}

\begin{frame}[fragile]{Window functions}
\begin{itemize}
\item Ak vynecháme PARTITION BY, partícia obsahuje všetky riadky výstupu.
\begin{alltt}
    SELECT salary, sum(salary) \alert{OVER ()}
    FROM employee;
\end{alltt}
\vspace*{4mm}
\pause
\item Okno možno pomenovať a použiť viackrát.
\begin{alltt}
    SELECT
        sum(salary) OVER \blue{w},
        avg(salary) OVER \blue{w}
    FROM employee
    \alert{WINDOW \blue{w} AS (PARTITION BY dept_id
                ORDER BY salary DESC);
}
\end{alltt}
\end{itemize}
\end{frame}

\begin{frame}[fragile]{Window frame}
\small
\begin{itemize}
\item Každý riadok má svoje vlastné okno (window frame) --- zoznam riadkov, nad ktorými sa vykoná window function.
\item Okno závisí od použitej funkcie. Default window frame možno v rámci dotazu modifikovať.
\item Komplikovaná syntax. Nasledujúce dve formulácie sú ekvivalentné významovo, ale vraj nie z hľadiska efektívnosti výpočtu\dots
\begin{alltt}
\blue{ROWS UNBOUNDED PRECEDING EXCLUDE CURRENT ROW}
\green{ROWS BETWEEN UNBOUNDED PRECEDING AND 1 PRECEDING}
\end{alltt}
\end{itemize}
\end{frame}

\begin{frame}[fragile]{Window frame}
\begin{alltt}
SELECT salary, \green{sum}(salary) \alert{OVER (ORDER BY salary)}
FROM employee;
\end{alltt}
{\scriptsize
\begin{alltt}
\blue{
 salary  |   sum
-{}-{}-{}-{}-{}-{}-{}-{}-{}+-{}-{}-{}-{}-{}-{}
  800.00 |   800.00
  950.00 |  1750.00
 1100.00 |  2850.00
 1300.00 |  4150.00
 1500.00 |  7150.00   \green{<---- pribudlo 3000, nie 1500}
 1500.00 |  7150.00
 1750.00 |  8900.00
}
\end{alltt}
}
Okno: riadky v partícii po daný riadok plus riadky s rovnakou hodnotou.
\end{frame}

\begin{frame}[fragile]{Vzťah M:N --- pomocná tabuľka}
Zamestnanec môže pracovať na viacerých projektoch a na projekte môže pracovať viacero zamestnancov.
Takýto vzťah nemožno uložiť v jednej hodnote:
\begin{itemize}
\item v riadku zamestnanca by bol zoznam projektov (pole),
\item v riadku projektu by bol zoznam zamestnancov.
\end{itemize}
Preto vzťah ukladáme do pomocnej tabuľky \verb|emp_projects|, kde je jeden riadok pre každú dvojicu (zamestnanec, projekt):
\bigskip

\begin{minipage}[t]{0.4\linewidth}
\footnotesize
\begin{alltt}
emp_id | proj_id
-{}-{}-{}-{}-{}-{}-{}+-{}-{}-{}-{}-{}-{}-{}-{}
    15 |       1
    15 |       2
    17 |       1
\end{alltt}
\end{minipage}
\hfill
\begin{minipage}[t]{0.55\linewidth}
\footnotesize
Zamestnanec 15 pracuje na projektoch 1 a 2.
Na projekte 1 pracujú zamestnanci 15 a 17.
\end{minipage}
\end{frame}

\begin{frame}[fragile]{Vzťah M:N --- pomocná tabuľka}
\begin{alltt}
CREATE TABLE project (
   proj_id INTEGER PRIMARY KEY,
   name TEXT NOT NULL
);

CREATE TABLE emp_projects (
   emp_id INTEGER NOT NULL,
   proj_id INTEGER NOT NULL
);
\end{alltt}
Join prebieha len cez rovnosť jednoduchých stĺpcov, takže sa dá indexovať a optimalizátor ho vie spracovať bežnými algoritmami pre join.
Pri poli by tento join už nebol rovnosťou jednoduchých stĺpcov:
museli by sme hľadať hodnotu medzi prvkami poľa.
\end{frame}

\begin{frame}[fragile]{Kedy má pole zmysel}
Hoci zložené hodnoty vo všeobecnosti nepoužívame, pole sa hodí, keď ide o:
\begin{itemize}
\item hodnotu, ktorú prirodzene tvorí n-tica, napr. súradnice bodu (x, y, z),
\item hodnoty rovnakého druhu, ktorých počet sa mení, a ktoré by inak zaberali veľa stĺpcov väčšinou s \texttt{NULL}, napr. namiesto stĺpcov \verb|telefon1| až \verb|telefon10|,
\item podmnožinu konečného zoznamu hodnôt, ktorú čítame a meníme ako celok, napr. tagy knihy \verb|'fantasy'|, \verb|'horor'|, \verb|'humor'|, ktoré nemajú vzťahy k iným tabuľkám,
\item výstup, ktorý zostavujeme z dát uložených nekompozitne, napr. zoznam projektov každého zamestnanca.
\end{itemize}
\end{frame}

\begin{frame}[fragile]{Array datatype}
Databázy bežne podporujú polia (nie však SQLite).\\
Dátový typ v PostgreSQL:
\begin{alltt}
  CREATE TABLE x (
      y  integer[4],
      z  text[][]
  );
\end{alltt}
Zápis v súlade s SQL štandardom:
\begin{alltt}
  CREATE TABLE x (
      y  integer ARRAY[4],
      z  text ARRAY[][]
  );
\end{alltt}
Aktuálna implementácia veľkosť poľa ignoruje.
\end{frame}

\begin{frame}[fragile]{Arrays}
\small
Pole možno vytvoriť priamo v rámci dotazu.
\begin{alltt}
SELECT ARRAY[1,2] || ARRAY[3,4] AS a;
SELECT ARRAY[1, 2] || '\{3, 4\}';
\end{alltt}
Na prácu s poľom existuje množstvo funkcií, napr. \verb|array_prepend|, \verb|array_length|.
\begin{alltt}
SELECT array_position(
    ARRAY['mon','tue','wed','thu','fri','sat','sun'],
    'mon'
  );
SELECT * FROM x WHERE 47 = ANY (some_array);
SELECT * FROM x WHERE 47 = ALL (some_array);
\end{alltt}
\end{frame}

\begin{frame}[fragile]{Arrays}
\small
Pole tiež možno získať aplikovaním agregačnej funkcie.
\begin{alltt}
SELECT e.name, \alert{array_agg}(p.name)
FROM employee e
    JOIN emp_projects ep ON ep.emp_id = e.emp_id
    JOIN project p ON p.proj_id = ep.proj_id
GROUP BY e.emp_id, e.name;
\end{alltt}
{\color{Cerulean}
\begin{alltt}
     name     |         array_agg
-{}-{}-{}-{}-{}-{}-{}-{}-{}-{}-{}-{}-{}-{}+-{}-{}-{}-{}-{}-{}-{}-{}-{}-{}-{}-{}-{}-{}-{}-{}-{}-{}-{}-{}-{}-{}-{}-{}-{}-{}-{}-{}
 MoneyMagnet  | \{Enviro1,Enviro2,Nuclear3\}
 SalesManager | \{Enviro2,Nuclear3\}
 InsightNinja | \{Nuclear3\}
 MrFixIt      | \{Enviro1\}
 DealMaster   | \{Enviro1\}
\end{alltt}
}
Tento dotaz nezaručuje konkrétne poradie prvkov vo výsledných poliach.
\end{frame}


\begin{frame}[fragile]{JSON a jsonb}
Čiastočne štruktúrované dáta: každý záznam môže mať iné atribúty.\\
PostgreSQL má typ \alert{\texttt{json}} (text) a \alert{\texttt{jsonb}} (binárna reprezentácia, indexovateľná).
\bigskip

\begin{alltt}
ALTER TABLE employee ADD COLUMN contact jsonb;

INSERT INTO employee (emp_id, name, contact)
VALUES (11, 'King', '\{"phone": "555-0101",
  "email": "king@example.com",
  "languages": ["sk", "en"]\}');
\end{alltt}
\end{frame}

\begin{frame}[fragile]{JSON --- prístup k hodnotám}
\begin{itemize}
\item \alert{\texttt{->}} a \alert{\texttt{-{}>{}>}}: prístup ku kľúču
\item \verb|contact->'phone'| vracia \verb|jsonb|
\item \texttt{contact-{}>{}>'phone'} vracia text,
\item \verb|contact ? 'email'| overí, či kľúč existuje,
\item \texttt{contact->'languages'-{}>{}>0} je prvý prvok poľa.
\end{itemize}
\begin{alltt}
SELECT name FROM employee
WHERE contact \alert{@>} '\{"email": "king@example.com"\}';
\end{alltt}
\end{frame}

\begin{frame}[fragile]{Operátor @> (containment)}
\verb|a @> b| je pravda, ak \verb|jsonb| hodnota \verb|a| obsahuje všetko, čo je v \verb|b| (kľúče aj hodnoty, aj vnorene).
{\footnotesize
\begin{alltt}
'\{"a": 1, "b": 2\}' \alert{@>} '\{"a": 1\}'    ------ TRUE
'\{"a": 1\}' \alert{@>} '\{"a": 1, "b": 2\}'    ------ FALSE
'["sk", "en"]' \alert{@>} '["en"]'          ------ TRUE
'\{"c": \{"ip": 1, "x": 2\}\}' \alert{@>} '\{"c": \{"ip": 1\}\}' ------ TRUE
'\{"c": \{"ip": 1\}\}' \alert{@>} '\{"ip": 1\}'   ------ FALSE
\end{alltt}
}
Pole sa porovnáva po prvkoch, na ich poradí nezáleží; vnorený objekt sa porovnáva rekurzívne.
\end{frame}

\begin{frame}[fragile]{Príklad: logovanie pre web (REST API)}
{}
{\scriptsize
\begin{alltt}
CREATE TABLE request_log (
  id SERIAL PRIMARY KEY,
  payload \alert{jsonb}
);

INSERT INTO request_log (payload) VALUES
 ('\{"method": "GET", "path": "/api/employees/11",
   "status": 200, "duration_ms": 12,
   "client": \{"ip": "10.0.0.7"\}\}'),
 ('\{"method": "GET", "path": "/api/projects",
   "status": 500, "duration_ms": 903,
   "client": \{"ip": "10.0.0.7"\}\}');

SELECT payload-{}>{}>'path', (payload-{}>{}>'duration\_ms')::int
FROM request_log
WHERE (payload-{}>{}>'status')::int >= 500;

SELECT id, payload->'client'-{}>{}>'ip'
FROM request_log
WHERE payload \alert{@>} '\{"client": \{"ip": "10.0.0.7"\}\}';
\end{alltt}
}
\end{frame}

\begin{frame}[fragile]{json\_extract --- JSON v SQLite}
% https://www.sqlite.org/json1.html#jsonb
SQLite nemá samostatný JSON dátový typ. JSON ukladáme ako text alebo (od verzie 3.45) vo formáte JSONB typu BLOB.
Prístup k hodnotám: \alert{\texttt{json\_extract}}, \alert{\texttt{->}}, \alert{\texttt{-{}>{}>}}.
{\footnotesize
\begin{alltt}
SELECT name, \alert{json_extract}(contact, '\$.phone') AS phone
FROM employee;

SELECT name FROM employee
WHERE \alert{json_extract}(contact, '\$.phone') = '555-0101';

SELECT name, \alert{json_extract}(contact, '\$.languages[0]')
FROM employee
WHERE contact IS NOT NULL;
\end{alltt}
}
{\footnotesize
\begin{alltt}
SELECT name,
       contact -> '\$.phone' AS json_val,
       contact \alert{-{}>{}>} '\$.phone' AS sql_val
FROM employee
WHERE contact IS NOT NULL;
\end{alltt}
}
\end{frame}

\begin{frame}[fragile]{Pole, jsonb, alebo tabuľka?}
\begin{itemize}
\item pole: krátky homogénny zoznam, ktorý čítame a meníme ako celok
\item jsonb: rôznorodé voliteľné atribúty bez vzťahov k iným tabuľkám, namiesto mnohých stĺpcov s väčšinou NULL
\item samostatná tabuľka: vzťahy s inými tabuľkami, cudzie kľúče, many-to-many relationships
\end{itemize}
\end{frame}

\begin{frame}{Literatúra}
\begin{itemize}
\item {\scriptsize\url{https://www.sqltutorial.org/sql-window-functions/}}
\item {\scriptsize\url{https://www.sqlite.org/windowfunctions.html}}
\item {\scriptsize\url{https://www.postgresql.org/docs/current/tutorial-window.html}}
\item {\scriptsize\url{https://www.postgresql.org/docs/current/arrays.html}}
\end{itemize}
\end{frame}

\end{document}
